\pdfvariable trailerid {[<C3112026090300000000000000000000><C3112026090300000000000000000000>]}
\documentclass[10pt]{article}
\usepackage[margin=0.72in]{geometry}
\usepackage{amsmath,amssymb,amsthm,booktabs,microtype,hyperref}
\hypersetup{colorlinks=true,linkcolor=blue,urlcolor=blue}
\providecommand{\CRevisionRound}{2}
\newtheorem{theorem}{Theorem}
\newcommand{\Rea}{\operatorname{Re}}
\newcommand{\tr}{\operatorname{tr}}
\setlength{\emergencystretch}{3em}
\title{A Convention-Explicit Hopf Atlas for the Brusselator}
\author{Route-A source-local certificate HCS-C311}
\date{3 September 2026}
\begin{document}\maketitle
\begin{abstract}
For every $A>0$ and $B\ge0$, we prove positive-quadrant global existence for
the Brusselator and give its complete node, focus, defective, and Hopf
equilibrium atlas.
\ifnum\CRevisionRound>0
With explicitly normalized right/left eigenvectors we derive the full complex
$G_{21}$, distinguish three coefficient conventions, prove uniform
supercriticality, and obtain leading cycle amplitude and frequency.
\fi
\ifnum\CRevisionRound>1
Independent exact, symbolic, replay, mutation, and deterministic-build lanes
audit the theorem and its strictly weak Route-A orbit signal.
\fi
\end{abstract}

\section{Global existence and complete linear atlas}
Consider
\begin{equation}\label{eq:model}
 \dot x=A-(B+1)x+x^2y,\qquad \dot y=Bx-x^2y,
 \quad A>0,\ B\ge0.
\end{equation}
At $x=0$ the first component equals $A>0$, and at $y=0$ the second is
$Bx\ge0$; hence the closed nonnegative quadrant is invariant.  There
$(x+y)'=A-x\le A$, so neither coordinate can blow up in finite time.
Every nonnegative solution is therefore global.

\begin{theorem}[All-parameter equilibrium chambers]\label{thm:linear}
Equation~\eqref{eq:model} has the unique nonnegative equilibrium
$E=(A,B/A)$.  Put $\tau=B-1-A^2$.  Its Jacobian is
\begin{equation}\label{eq:J}
 J=\begin{pmatrix}B-1&A^2\\-B&-A^2\end{pmatrix},
 \qquad \tr J=\tau,\quad\det J=A^2.
\end{equation}
Within $B\ge0$, the increasing-$B$ atlas is
\begin{center}\small
\begin{tabular}{@{}ll@{}}\toprule
range or boundary & type \\ \midrule
$B<(A-1)^2$ & stable node \\
$B=(A-1)^2$ & defective stable node \\
$(A-1)^2<B<1+A^2$ & stable focus \\
$B=1+A^2$ & Hopf pair $\pm iA$ \\
$1+A^2<B<(A+1)^2$ & unstable focus \\
$B=(A+1)^2$ & defective unstable node \\
$B>(A+1)^2$ & unstable node \\ \bottomrule
\end{tabular}
\end{center}
Empty ranges are omitted; in particular $A=1,B=0$ is the stable defective
boundary.
\end{theorem}
\begin{proof}
At equilibrium, adding the equations gives $x=A$, then the second gives
$y=B/A$.  Formula~\eqref{eq:J} follows by differentiation.  Its eigenvalues
are $\tau/2\pm\sqrt{\tau^2-4A^2}/2$.  The signs of $\tau$ and of this
discriminant give every row.  At equality the positive off-diagonal entry
$A^2$ prevents $J$ from being scalar, hence the repeated eigenvalue is
defective.
\end{proof}

\ifnum\CRevisionRound>0
\section{Exact normalized Hopf coefficient}
Set $B=1+A^2$ and translate $u=x-A$, $v=y-B/A$.  The nonlinear expansion is
\begin{align*}
 \dot u&=A^2u+A^2v+(A+A^{-1})u^2+2Auv+u^2v,\\
 \dot v&=-(A^2+1)u-A^2v-(A+A^{-1})u^2-2Auv-u^2v.
\end{align*}
Fix right and adjoint eigenvectors
\begin{equation}\label{eq:qp}
 q=(1,-1+i/A)^{\mathsf T},\qquad
 p=((1+iA)/2,iA/2)^{\mathsf T},\qquad \bar p^{\mathsf T}q=1.
\end{equation}
If $\mathcal B,\mathcal C$ are the second and third derivative multilinear
forms, use the convention
\begin{align}\label{eq:G}
G_{21}=\bar p^{\mathsf T}\{&\mathcal C(q,q,\bar q)
-2\mathcal B(q,J^{-1}\mathcal B(q,\bar q))\\
&+\mathcal B(\bar q,(2iA I-J)^{-1}\mathcal B(q,q))\}.
\end{align}
Direct substitution, with no numerical fitting, gives
\begin{equation}\label{eq:Gvalue}
 G_{21}=-\left(1+\frac2{A^2}\right)
 -i\frac{4A^4-7A^2+4}{3A^3},\qquad
 \ell_1=\frac{\Rea G_{21}}{2A}=-\frac{A^2+2}{2A^3}<0.
\end{equation}

\begin{theorem}[Uniformly supercritical branch]\label{thm:hopf}
Let $\mu=B-(1+A^2)$.  The Hopf crossing is transverse with real-part
derivative $1/2$.  For all sufficiently small $\mu>0$ there is one local
periodic-orbit branch modulo phase; it is stable.  In the coordinate
of~\eqref{eq:qp},
\begin{align}
 r^2&=\frac{A^2}{A^2+2}\mu+O(\mu^2),\label{eq:r}\\
 \Omega&=A-\frac{4A^4-7A^2+4}{6A(A^2+2)}\mu+O(\mu^2).\label{eq:freq}
\end{align}
The first harmonic of $x-A$ has amplitude
$2A\sqrt{\mu/(A^2+2)}$ up to an $O(\mu)$ coordinate correction.
\end{theorem}
\begin{proof}
Near zero the eigenvalues are
$\mu/2\pm i\sqrt{A^2-\mu^2/4}$, proving transversality and absence of linear
frequency detuning.  The physical-time normal form is
$\dot w=(\mu/2+iA)w+(G_{21}/2)w|w|^2+$ higher terms.  Its radial cubic
coefficient is $-(A^2+2)/(2A^2)$, whereas $\ell_1$ in~\eqref{eq:Gvalue}
includes the additional division by $A$.  Solving the radial equation gives
\eqref{eq:r}; its derivative is negative at the nonzero root.  The imaginary
part of $G_{21}/2$ gives~\eqref{eq:freq}, and $u=w+\bar w+O(|w|^2)$ gives the
first harmonic.
\end{proof}

The face $B=0$ is included in Theorem~\ref{thm:linear}.  The face $A=0$ is
not: $B/A$ is singular and $x=0$ becomes an equilibrium line.  No global
uniqueness of periodic orbits is inferred from the local Hopf theorem.
\fi

\ifnum\CRevisionRound>1
\section{Evidence, collisions, and Route-A boundary}
Twelve rational values $1/4\le A\le10$ produce 55 admissible chamber probes
and 827 audited leaves.  An independent checker performs 742 exact/numerical
checks.  A separate SymPy program differentiates the vector field and derives
\eqref{eq:Gvalue} in 13 identities; replay is byte exact and 26 repaired-hash
or parser attacks must fail.  These lanes audit conventions, not the proof.

C249 owns the global Van der Pol Li\'enard cycle, C235 a simplex replicator,
and C254 a washout/transcritical chemostat.  C311 instead closes the
Brusselator's all-$A$ chemical Hopf coefficient and linear chamber geometry.

The strict tuple is
\[
(\mathrm{A0\_FAIL},\mathrm{A1\_WEAK},\mathrm{A2\_FAIL},
 \mathrm{A3\_FAIL},\mathrm{A4\_FAIL}).
\]
The local isolated cycle is genuine source dynamics, earning only A1 weak.
Continuous reaction parameters carry no rational-prime local owner (A0) or
logarithmic-prime clock (A2); the normal form gives no target determinant
(A3) or self-adjoint target-zero operator (A4).  Route A is rejected and
Route B remains locked under \texttt{NO\_BAD\_EULER\_OR\_ROOT\_NUMBER}.
No target Euler factor, root number, automorphy statement, divisor law,
functional equation, or zero correspondence is claimed.

\paragraph{AI use.} A generative language model assisted drafting and code
scaffolding.  The displayed derivation, independent symbolic reconstruction,
adversarial tests, and archived artifacts define the audit record.
\fi

\section*{Source lineage}
\begin{thebibliography}{9}
\bibitem{PL1968} I. Prigogine and R. Lefever,
``Symmetry breaking instabilities in dissipative systems. II,''
\emph{J. Chem. Phys.} 48 (1968), 1695--1700.
DOI: \href{https://doi.org/10.1063/1.1668896}{10.1063/1.1668896}.
\bibitem{Kuznetsov} Y. A. Kuznetsov,
\emph{Elements of Applied Bifurcation Theory}, 3rd ed., Springer, 2004.
DOI: \href{https://doi.org/10.1007/978-1-4757-3978-7}{10.1007/978-1-4757-3978-7}.
\end{thebibliography}
\end{document}
